% !TEX TS-program = xelatex
\documentclass[12pt,a4paper]{report}

% -- Math (must come before bidi/polyglossia) ----------------
\usepackage{amsmath, amssymb}

% -- Graphics & floats (must come before bidi) ---------------
\usepackage{graphicx}
\usepackage{float}
\usepackage{caption}

% -- Tables --------------------------------------------------
\usepackage{booktabs}
\usepackage{multirow}
\usepackage{array}
\usepackage{tabularx}

% -- Colors & listings (must come before bidi) ---------------
\usepackage{xcolor}
\usepackage{listings}
\lstset{
  basicstyle=\ttfamily\small,
  breaklines=true,
  frame=single,
  backgroundcolor=\color{gray!8}
}

% -- TikZ (must come before bidi) ----------------------------
\usepackage{tikz}
\usetikzlibrary{shapes.geometric, arrows.meta, positioning, fit, calc, backgrounds}

% -- Hyperlinks (before bidi) --------------------------------
\usepackage{hyperref}
\hypersetup{colorlinks=true, linkcolor=black, citecolor=blue!60!black}

% -- Enumitem ------------------------------------------------
\usepackage{enumitem}

% -- References (before bidi) --------------------------------
\usepackage[numbers]{natbib}
\bibliographystyle{plainnat}

% -- Font and language (XeLaTeX) - MUST come last ------------
\usepackage{fontspec}
\usepackage{polyglossia}
\setmainlanguage{english}
\setotherlanguage{arabic}
\newfontfamily\arabicfont[Script=Arabic]{Amiri}

% -- Custom colours for TikZ figures -------------------------
\definecolor{pipeblue}{RGB}{52, 120, 190}
\definecolor{pipegreen}{RGB}{56, 142, 60}
\definecolor{pipeorange}{RGB}{230, 120, 30}
\definecolor{pipegray}{RGB}{110, 110, 110}
\definecolor{pipeyellow}{RGB}{180, 140, 0}
\definecolor{pipered}{RGB}{192, 57, 43}
\definecolor{pipepurple}{RGB}{120, 60, 160}

% -- Inline bibliography for one-file testing ----------------
\begin{filecontents*}{bibliography.bib}
@misc{retasyDataset,
  title        = {RetaSy Quranic Audio Dataset},
  howpublished = {\url{https://huggingface.co/datasets/RetaSy/quranic_audio_dataset}},
  note         = {Dataset used for learner-oriented Quranic recitation experiments}
}

@misc{quranJsonTajwid,
  title        = {Quran JSON Tajwid Reference},
  note         = {Structured Quran and Tajweed rule reference used to derive expected rule positions}
}

@misc{quranMd,
  title        = {Quran-MD Audio Dataset},
  note         = {Clean Quran recitation data explored for scalability, word-level, and ayah-level experiments}
}
\end{filecontents*}

% -- Title ---------------------------------------------------
\title{\textbf{A Modular Architecture for Automatic Tajweed Assessment\\
               of Quranic Recitation}}
\author{Test Compilation File}
\date{\today}

% ============================================================
\begin{document}

\maketitle
\tableofcontents
\chapter{Development, Experimentation, and Results}
\label{chap:development_experimentation}

This chapter presents the development state of the proposed Tajweed assessment
system, the main techniques used in the implementation, and the experimental
results obtained during the project. The goal is not only to report the final
numbers, but also to explain how the system evolved from the first prototype to
the current modular architecture.

The system was designed as a modular pipeline instead of a single global model.
This decision was important because Tajweed errors do not all have the same
nature. Some errors are related to duration, such as \textit{madd} and
\textit{ghunnah}. Some errors are related to transitions between letters, such
as \textit{ikhfa} and \textit{idgham}. Others are related to articulation and
burst energy, such as \textit{qalqalah}. Finally, content errors concern whether
the learner recited the expected words or not. A single model would make all
these problems mixed together, while a modular system makes each problem easier
to train, evaluate, debug, and explain.

\section{Development Evolution}

The project evolved through several stages. Each stage was motivated by an
observed limitation in the previous version.

\begin{table}[H]
\centering
\small
\begin{tabularx}{\textwidth}{p{0.22\textwidth} X X}
\toprule
\textbf{Stage} & \textbf{Work done} & \textbf{Main observation} \\
\midrule
Initial setup &
Created the first manifests, audio loading scripts, feature extraction, and
training scripts. &
The project needed a stable data format before model training could become
reliable. \\
\midrule
First modular version &
Separated the system into duration, transition, burst, and content modules. &
Specialist modules were easier to evaluate than a single monolithic classifier. \\
\midrule
Full-verse content attempt &
Trained content recognition on complete verses. &
Full-verse recognition was too difficult with the available learner data. \\
\midrule
Chunked content reformulation &
Split content into shorter chunks and trained a chunked content model. &
Shorter sequences improved content recognition strongly. \\
\midrule
Hardcase mining &
Analyzed confusion matrices and re-used difficult examples for training. &
The system improved when training focused on repeated errors. \\
\midrule
Localized support models &
Added models that detect acoustic evidence around expected rule positions. &
This improved interpretability because the system could inspect local evidence,
not only whole-clip labels. \\
\midrule
Fusion and threshold tuning &
Combined conservative predictions with localized evidence and tuned thresholds. &
Promotion was accepted only when side-by-side evaluation improved the baseline. \\
\midrule
Scalability experiments &
Explored Quran-MD word and ayah data, weak routing labels, and larger
vocabularies. &
Large clean recitation data helps scalability, but does not replace learner
mistake annotations. \\
\bottomrule
\end{tabularx}
\caption{Development evolution from the first prototype to the current system.}
\label{tab:development_evolution}
\end{table}

\section{Datasets and Textual References}

The system uses JSONL manifest files as the central corpus format. Each line in
a manifest represents one sample or one segment. A row can contain the audio
path, normalized text, surah and ayah information, reciter metadata, rule labels,
rule positions, and routing targets.

\subsection{RetaSy Learner Recitation Dataset}

The main learner-oriented dataset used in the project is the RetaSy Quranic
audio dataset. This dataset is important because it contains recitations from
learners, not only clean professional reciters. For a Tajweed assessment system,
learner data is more useful than perfect recitation alone because the model must
learn how mistakes appear acoustically.

However, the dataset is still limited for this task. The number of distinct
verses and annotated learner errors is small compared with the full Quran, which
contains more than six thousand verses. This limitation explains why the system
promotes conservative baselines and treats larger-data experiments as research
extensions until they are validated.

\subsection{Quran JSON and Colored Mushaf Reference}

The project also used a structured Quran reference based on Quran JSON and the
colored Mushaf idea. In the colored Mushaf, Tajweed rules are visually marked by
colors. In the project, this information was used in structured form through the
local reference files under \texttt{external/quranjson-tajwid}.

This reference was useful because audio data alone does not explicitly tell the
system where a Tajweed rule should appear. The Quran JSON reference provides the
canonical Quran text, surah and ayah identifiers, and rule-related spans. These
spans were transformed into project manifests such as
\texttt{quranjson\_rules.jsonl} and joined with learner audio manifests such as
\texttt{retasy\_quranjson\_train.jsonl}.

In simple words, the Quran JSON and colored Mushaf reference tells the system
what should exist in the verse, while the audio model tries to detect whether
the learner produced it correctly.

\subsection{Quran-MD Clean Recitation Data}

The Quran-MD dataset was explored as a larger source of clean Quran recitation.
It was useful for word-level and ayah-level scalability experiments because it
contains many more audio clips than the learner-oriented subsets. The main
benefit of this dataset is coverage: it exposes the content model and routing
experiments to more Arabic words and more Quranic text.

The limitation is that Quran-MD is mainly clean recitation data. Therefore, it
cannot directly solve learner mistake detection. It is useful for pretraining,
representation learning, and scalability, but it must be validated against
learner data before being promoted.

\section{Corpus Construction}

The corpus construction step transforms raw audio metadata and textual
references into model-ready manifests. The simplified logic is shown in
Listing~\ref{lst:corpus_pseudocode}.

\begin{lstlisting}[caption={Simplified pseudocode for corpus construction.},label={lst:corpus_pseudocode}]
input:
    audio files
    metadata
    Quran JSON reference text
    colored Mushaf Tajweed rule reference
    normalization rules

for each sample:
    load metadata
    normalize Arabic text
    map sample to surah and ayah
    join sample with Quran JSON verse
    extract expected Tajweed rule positions

    assign routing labels:
        use_duration
        use_transition
        use_burst
        use_content

    if the row contains madd or ghunnah:
        add to duration manifest

    if the row contains ikhfa or idgham:
        add to transition manifest

    if the row contains qalqalah:
        add to burst manifest

    if expected text is available:
        create full verse or chunked content targets

    write JSONL row with:
        sample id
        audio path
        normalized text
        rule labels
        rule positions
        routing labels
        source dataset
        reciter metadata
\end{lstlisting}

\section{Feature Representation}

\subsection{MFCC Features}

The rule-based modules use Mel-Frequency Cepstral Coefficients (MFCCs). MFCCs
are compact acoustic features that describe the spectral shape of speech. They
are useful for local Tajweed phenomena because rules such as duration,
transitions, and qalqalah depend on short acoustic events.

The system also uses delta and delta-delta features. A delta feature describes
how the signal changes over time, while a delta-delta feature describes the
change of that change. This helps the model capture movement in speech, not
only static sound frames.

\subsection{Wav2Vec-Style SSL Features}

For content recognition, the system uses self-supervised speech representations
inspired by Wav2Vec-style models. Self-supervised learning (SSL) means that a
large speech model learns useful audio representations from audio itself before
being adapted to a specific task.

Content recognition is closer to automatic speech recognition than simple rule
classification. For that reason, stronger learned speech representations are
more suitable than handcrafted MFCCs. The content model uses these features and
then predicts Arabic characters.

\section{Modeling Techniques}

\subsection{BiLSTM Temporal Modeling}

Several modules use a bidirectional Long Short-Term Memory network (BiLSTM). A
BiLSTM reads an audio sequence from left to right and from right to left. This
is useful because a Tajweed rule can depend on what comes before and after a
letter. For example, transition rules depend on the relationship between two
neighboring letters.

\subsection{CTC Loss}

Connectionist Temporal Classification (CTC) is used for content recognition.
CTC is useful when the training data contains an audio clip and its text, but
does not contain exact timestamps for every character. Instead of requiring a
manual alignment between frames and letters, CTC learns the alignment during
training.

The decoder converts the CTC frame probabilities into a final text prediction.
The project tested several decoding choices, including greedy decoding,
open decoding, lexicon-constrained decoding, and blank-penalty tuning.

\subsection{Hardcase Mining}

Hardcase mining is a targeted error-analysis technique. Instead of only looking
at the average accuracy, the system extracts examples where the model makes
confident mistakes. These examples are then used to guide retraining or
comparison experiments.

\begin{lstlisting}[caption={Hardcase mining loop.},label={lst:hardcase_loop}]
train baseline model
evaluate model
extract high-confidence mistakes
group errors by rule and confusion type
build hardcase list
retrain or tune with hardcase awareness
compare against the previous baseline
\end{lstlisting}

\subsection{Localized Support Models}

A localized support model tries to find acoustic evidence near the expected
position of a Tajweed rule. This is different from a whole-verse classifier.
The whole-verse classifier predicts one label for the entire clip, while the
localized support model checks whether the evidence exists around a specific
time region.

This improves explainability. When the system reports a rule error, it can also
store whether localized evidence supported or contradicted the decision.

\subsection{Learned Fusion}

Fusion means combining several sources of evidence. In the duration path, the
system combines the conservative duration prediction with localized evidence.
This learned fusion was treated as experimental first. It was promoted only
after a strict verse-held-out evaluation showed improvement over the
conservative baseline.

\subsection{Threshold and Decoder Tuning}

Many models output probabilities. A threshold decides when a probability becomes
a positive prediction. Threshold tuning was used to improve rule decisions and
avoid blindly accepting low-confidence predictions.

For CTC content recognition, the blank symbol can dominate the output. If the
decoder predicts too many blanks, it deletes many characters. The project tuned
a blank penalty to reduce deletion errors and improve content predictions.

\subsection{Severity-Aware Error Weighting}

The final feedback does not treat all mistakes equally. A critical content error
or a strong rule violation should have more pedagogical weight than a minor
uncertainty. The project therefore uses severity-aware scoring.

This is close to an error-priority mechanism. It is not a neural attention
module in the deep-learning sense, but it gives more importance to errors that
matter more for feedback.

\begin{equation}
\text{weighted penalty}
= \text{error count} \times \text{severity weight} \times \text{confidence}.
\end{equation}

\section{System Architecture}

The system follows a route-then-specialize architecture. A routing layer decides
which specialist module should be used for a sample. The specialist modules then
produce predictions, and the feedback layer converts those predictions into
human-readable comments.

\begin{figure}[H]
\centering
\begin{tikzpicture}[
    node distance=1.1cm,
    block/.style={rectangle, rounded corners, draw=black!50, thick,
                  align=center, minimum height=0.8cm, minimum width=2.6cm,
                  fill=gray!8},
    route/.style={rectangle, rounded corners, draw=pipeblue, thick,
                  align=center, minimum height=0.8cm, minimum width=2.8cm,
                  fill=pipeblue!10},
    module/.style={rectangle, rounded corners, draw=pipegreen!70!black, thick,
                   align=center, minimum height=0.8cm, minimum width=2.7cm,
                   fill=pipegreen!10},
    output/.style={rectangle, rounded corners, draw=pipered!70!black, thick,
                   align=center, minimum height=0.8cm, minimum width=2.8cm,
                   fill=pipered!10},
    arrow/.style={-{Latex[length=2.2mm]}, thick}
]

\node[block] (audio) {Audio sample};
\node[block, below=of audio] (features) {Feature extraction\\MFCC or SSL};
\node[route, below=of features] (router) {Routing layer};

\node[module, below left=1.0cm and 3.2cm of router] (duration) {Duration\\madd, ghunnah};
\node[module, below left=1.0cm and 0.8cm of router] (transition) {Transition\\ikhfa, idgham};
\node[module, below right=1.0cm and 0.8cm of router] (burst) {Burst\\qalqalah};
\node[module, below right=1.0cm and 3.2cm of router] (content) {Content\\text correctness};

\node[output, below=2.0cm of router] (feedback) {Diagnosis and feedback};

\draw[arrow] (audio) -- (features);
\draw[arrow] (features) -- (router);
\draw[arrow] (router) -- (duration);
\draw[arrow] (router) -- (transition);
\draw[arrow] (router) -- (burst);
\draw[arrow] (router) -- (content);
\draw[arrow] (duration) -- (feedback);
\draw[arrow] (transition) -- (feedback);
\draw[arrow] (burst) -- (feedback);
\draw[arrow] (content) -- (feedback);

\end{tikzpicture}
\caption{General modular architecture of the Tajweed assessment system.}
\label{fig:modular_architecture}
\end{figure}

\section{Module Development}

\subsection{Duration Module}

The duration module detects rules related to length and nasal duration,
especially \textit{madd} and \textit{ghunnah}. It was built in several layers.

\paragraph{Dataset and labels.}
The duration corpus was created from samples that contain expected duration
rules. The Quran JSON and colored Mushaf reference helped identify where
\textit{madd} and \textit{ghunnah} should occur in the text. The manifest stores
the audio path, normalized text, expected rule, and target position.

\paragraph{MFCC feature pipeline.}
The audio is converted into MFCC features with temporal derivatives. This gives
the model a compact representation of the sound and its changes over time.

\paragraph{Sequence duration model.}
The main duration classifier analyzes the audio sequence and predicts whether
the target rule is correctly represented. Earlier duration experiments reached
approximately 0.838 focus accuracy in confusion analysis.

\paragraph{Localized support model.}
A second model was trained to detect local acoustic support around expected
duration positions. This improves the explanation because the system can inspect
the evidence near the expected rule location.

\paragraph{Learned fusion layer.}
The duration path later combined the conservative duration classifier with
localized evidence. This learned fusion was not accepted immediately. It was
tested against the conservative baseline on a strict verse-held-out subset.

\paragraph{Strict held-out approval.}
The learned fusion was promoted because it improved strict verse-held-out
accuracy from 0.954 to 0.973. The \textit{ghunnah} class improved from 0.745 to
0.872, and \textit{madd} improved from 0.980 to 0.985.

\subsection{Transition Module}

The transition module detects rules that happen between letters, mainly
\textit{ikhfa} and \textit{idgham}. In the official baseline, the first version
used one dominant transition label per verse or clip. This simplification made
training stable and allowed a clear baseline.

However, one verse can contain more than one rule family. For example, a verse
may contain both an \textit{ikhfa} case and another transition behavior. For
this reason, a multi-label transition experiment was also developed. In the
multi-label setting, the model can predict more than one transition label for a
sample. This version is promising, but the single-label hardcase transition
model remains the official promoted baseline because it currently gives the
strongest validated result.

The transition module improved from about 0.770 accuracy in early evaluations
to 0.901 official accuracy after hardcase-aware training and threshold tuning.

\subsection{Burst Module}

The burst module focuses on \textit{qalqalah}. Qalqalah is an articulation rule
where a specific consonant has a bouncing or echo-like burst. This acoustic
phenomenon is different from duration or transition, so it was treated as a
separate specialist module.

The promoted burst baseline reached 0.874 accuracy. The class-level analysis
showed 0.914 accuracy for the \textit{none} class and 0.814 accuracy for
\textit{qalqalah}. This means the module is usable, but qalqalah detection still
has room for improvement.

\subsection{Content Module}

The content module checks whether the recited text matches the expected text.
This was the most difficult part of the project because it is close to Arabic
automatic speech recognition.

The first full-verse content model performed poorly, with exact match around
0.016 and character accuracy around 0.536. This showed that full-verse
recognition was too hard for the available data and model size.

The project then reformulated content recognition as a chunked task. Instead of
asking the model to recognize a complete verse, the verse is split into shorter
segments. This improved the official chunked content baseline to 0.700 exact
match and 0.892 character accuracy.

Additional content experiments were tested:

\begin{itemize}[leftmargin=*]
    \item A 5000-word Quran-MD auxiliary model improved word-level recognition
    on clean held-out words, but direct transfer to chunked learner content was
    not reliable.
    \item A tiny auxiliary fine-tuning stage slightly improved the chunked
    baseline from 0.700 to 0.707 exact match, but it remains experimental.
    \item Full-ayah content experiments improved character-level behavior but
    were not reliable enough for learner-level open recognition.
\end{itemize}

\subsection{Routing Module}

The routing module decides which specialist module should be activated. The
official system uses rule and metadata information from the manifest. For
example, if a sample contains an expected \textit{madd} or \textit{ghunnah}
target, the duration module is activated. If it contains an \textit{ikhfa} or
\textit{idgham} target, the transition module is activated.

A learned router was also explored. It achieved strong weak-label consistency
on a mixed Retasy and Quran-MD routing dataset, but weaker Retasy-only
validation showed that it should remain experimental. Therefore, the official
baseline keeps the conservative rule-based router.

\section{Official Results}

Table~\ref{tab:official_results} summarizes the official promoted results.

\begin{table}[H]
\centering
\small
\begin{tabularx}{\textwidth}{p{0.28\textwidth} p{0.32\textwidth} X}
\toprule
\textbf{Module} & \textbf{Metric} & \textbf{Result} \\
\midrule
Duration & Suite accuracy & 0.993 \\
Duration & Strict verse-held-out accuracy & 0.973 \\
Duration: ghunnah & Suite class accuracy & 0.984 \\
Duration: madd & Suite class accuracy & 0.995 \\
Transition & Accuracy & 0.901 \\
Transition: none & Class accuracy & 0.896 \\
Transition: ikhfa & Class accuracy & 0.921 \\
Transition: idgham & Class accuracy & 0.857 \\
Burst / qalqalah & Accuracy & 0.874 \\
Content chunked baseline & Exact match & 0.700 \\
Content chunked baseline & Character accuracy & 0.892 \\
Content chunked baseline & Average edit distance & 0.731 \\
\bottomrule
\end{tabularx}
\caption{Official promoted results of the modular system.}
\label{tab:official_results}
\end{table}

\section{Ablation Study}

An ablation study measures the value added by each component. The principle is
simple: remove or replace one component, then compare performance. This section
reports both system-level and module-level ablations. In this project, the
ablation study was not only theoretical. Each important addition was evaluated
against the previous baseline before it was accepted or kept experimental.

\subsection{Added Components and Their Measured Effect}

Table~\ref{tab:added_components_ablation} summarizes the main components added
during development and the measurable effect of each one. The most important
rule used during development was conservative promotion: a new component was
accepted only when it improved the relevant validation result without creating a
hidden regression.

\begin{table}[H]
\centering
\small
\begin{tabularx}{\textwidth}{p{0.26\textwidth} p{0.23\textwidth} p{0.23\textwidth} X}
\toprule
\textbf{Added component} & \textbf{Before} & \textbf{After} & \textbf{Decision} \\
\midrule
Chunked content reformulation &
Full-verse exact 0.016, character accuracy 0.536 &
Chunked exact 0.700, character accuracy 0.892 &
Promoted. This was the largest content gain because short chunks are much
easier than complete verses. \\
\midrule
Content decoder blank penalty &
Blank penalty 1.0: exact 0.691, character accuracy 0.890, edit distance 0.741 &
Blank penalty 1.6: exact 0.700, character accuracy 0.892, edit distance 0.731 &
Promoted. The decoder deleted fewer characters and improved exact match. \\
\midrule
Duration learned fusion &
Conservative strict held-out accuracy 0.954; ghunnah 0.745; madd 0.980 &
Learned fusion strict held-out accuracy 0.973; ghunnah 0.872; madd 0.985 &
Promoted. It improved generalization on a verse-held-out subset. \\
\midrule
Transition hardcase training &
Early transition accuracy about 0.770 &
Official transition accuracy 0.901; ikhfa 0.921; idgham 0.857 &
Promoted. Hardcase training reduced repeated confusion between transition
classes. \\
\midrule
Dedicated burst module &
No separate qalqalah specialist &
Burst accuracy 0.874; qalqalah class accuracy 0.814 &
Kept as baseline. It adds a separate articulation pathway. \\
\midrule
Tiny auxiliary content fine-tuning &
Official chunked exact 0.700; character accuracy 0.892; critical errors 125 &
Tiny auxiliary v2 exact 0.707; character accuracy 0.893; critical errors 122 &
Kept as candidate. The gain is positive but small, so the original baseline is
preserved. \\
\midrule
Quran-MD word auxiliary model &
Old content model on word clips: exact 0.001, character accuracy 0.062 &
Word auxiliary model: exact 0.492, character accuracy 0.834 &
Useful for future pretraining. Direct transfer to chunked learner content failed,
so it was not promoted. \\
\midrule
Learned router with thresholds &
Default weak-label exact 0.895, macro-F1 0.976 &
Tuned thresholds exact 0.921, macro-F1 0.979 &
Experimental. Strong on weak labels, but Retasy-only validation was weaker, so
the rule-based router remains official. \\
\bottomrule
\end{tabularx}
\caption{Measured value added by the main components introduced during development.}
\label{tab:added_components_ablation}
\end{table}

\subsection{System-Level Ablation}

\begin{table}[H]
\centering
\small
\begin{tabularx}{\textwidth}{p{0.31\textwidth} p{0.24\textwidth} X}
\toprule
\textbf{System variant} & \textbf{Observed behavior} & \textbf{Conclusion} \\
\midrule
Single full-verse content direction &
Very low exact match on learner content. &
A monolithic content-first system is not sufficient. \\
\midrule
Specialist rule modules &
Duration, transition, and burst obtained stable class-level results. &
Modularity improves diagnosability and performance. \\
\midrule
Without localized evidence &
Predictions were less explainable and harder to validate. &
Localized support adds interpretability and safer fusion. \\
\midrule
Without hardcase mining &
Repeated high-confidence mistakes remained visible. &
Hardcase mining helps target real model weaknesses. \\
\midrule
Without severity weighting &
All errors had equal importance. &
Severity weighting makes feedback closer to teacher priorities. \\
\bottomrule
\end{tabularx}
\caption{System-level ablation summary.}
\label{tab:system_ablation}
\end{table}

\subsection{Duration Module Ablation}

\begin{table}[H]
\centering
\small
\begin{tabularx}{\textwidth}{p{0.30\textwidth} p{0.25\textwidth} X}
\toprule
\textbf{Variant} & \textbf{Result} & \textbf{Interpretation} \\
\midrule
Basic duration sequence model &
Focus accuracy around 0.838 in earlier analysis. &
The model learned useful duration patterns but confused \textit{ghunnah} and
\textit{madd}. \\
\midrule
Hardcase-aware duration training &
Improved some difficult cases but could reduce general balance if overused. &
Hardcase weighting must be controlled carefully. \\
\midrule
Conservative duration baseline &
Strict held-out accuracy 0.954. &
Reliable baseline before fusion. \\
\midrule
Learned duration fusion &
Strict held-out accuracy 0.973. &
Fusion improved the conservative baseline and was promoted. \\
\bottomrule
\end{tabularx}
\caption{Ablation inside the duration module.}
\label{tab:duration_ablation}
\end{table}

\subsection{Transition Module Ablation}

\begin{table}[H]
\centering
\small
\begin{tabularx}{\textwidth}{p{0.30\textwidth} p{0.25\textwidth} X}
\toprule
\textbf{Variant} & \textbf{Result} & \textbf{Interpretation} \\
\midrule
Early transition classifier &
Accuracy around 0.770. &
The initial model detected some transition patterns but still confused
\textit{none}, \textit{ikhfa}, and \textit{idgham}. \\
\midrule
Hardcase transition checkpoint &
Official accuracy 0.901. &
Hardcase-aware training significantly improved transition reliability. \\
\midrule
Localized transition support &
Provided local evidence for transition decisions. &
Useful for explanation and hybrid inference. \\
\midrule
Multi-label transition experiment &
Exact match around 0.819 in gold-only profile experiments. &
Promising for verses with multiple rules, but not promoted yet. \\
\bottomrule
\end{tabularx}
\caption{Ablation inside the transition module.}
\label{tab:transition_ablation}
\end{table}

\subsection{Burst Module Ablation}

\begin{table}[H]
\centering
\small
\begin{tabularx}{\textwidth}{p{0.30\textwidth} p{0.25\textwidth} X}
\toprule
\textbf{Variant} & \textbf{Result} & \textbf{Interpretation} \\
\midrule
No burst specialist &
Qalqalah would be mixed with other rule errors. &
The system would lose a dedicated articulation pathway. \\
\midrule
Dedicated burst classifier &
Accuracy 0.874. &
The specialist module provides a stable qalqalah baseline. \\
\midrule
Class-level analysis &
None 0.914, qalqalah 0.814. &
The module is stronger on non-qalqalah cases than positive qalqalah cases. \\
\bottomrule
\end{tabularx}
\caption{Ablation inside the burst module.}
\label{tab:burst_ablation}
\end{table}

\subsection{Content Module Ablation}

\begin{table}[H]
\centering
\small
\begin{tabularx}{\textwidth}{p{0.32\textwidth} p{0.25\textwidth} X}
\toprule
\textbf{Variant} & \textbf{Result} & \textbf{Interpretation} \\
\midrule
Full-verse CTC content &
Exact 0.016, character accuracy 0.536. &
Full verses are too difficult for the available learner data. \\
\midrule
First chunked content model &
Exact around 0.383. &
Chunking immediately made the task more learnable. \\
\midrule
Official chunked baseline &
Exact 0.700, character accuracy 0.892. &
This is the current promoted content baseline. \\
\midrule
Blank-penalty tuned decoder &
Improved deletion behavior. &
Decoder tuning matters because CTC can over-predict blanks. \\
\midrule
Quran-MD word auxiliary model &
Clean word exact 0.492, character accuracy 0.834. &
Useful for learning words, but direct transfer to learner chunks failed. \\
\midrule
Tiny auxiliary transfer &
Exact 0.707, character accuracy 0.893. &
Small gain, kept experimental until more validation. \\
\midrule
Full-ayah content experiment &
Exact 0.063, character accuracy 0.762. &
Promising for future scalability, not ready for official use. \\
\bottomrule
\end{tabularx}
\caption{Ablation inside the content module.}
\label{tab:content_ablation}
\end{table}

\subsection{Routing and Scoring Ablation}

\begin{table}[H]
\centering
\small
\begin{tabularx}{\textwidth}{p{0.31\textwidth} p{0.24\textwidth} X}
\toprule
\textbf{Variant} & \textbf{Result} & \textbf{Interpretation} \\
\midrule
Rule-based router &
Official routing path. &
Reliable because it uses known rule metadata. \\
\midrule
Learned router v5 &
Weak-label exact around 0.921 on mixed data, weaker on Retasy-only validation. &
Promising but not official because weak labels are not enough. \\
\midrule
Unweighted scoring &
All errors have equal cost. &
Pedagogically weak because serious mistakes are not prioritized. \\
\midrule
Severity-aware scoring &
Critical errors receive higher penalty. &
More useful for learner feedback and teacher interpretation. \\
\bottomrule
\end{tabularx}
\caption{Ablation inside routing and scoring.}
\label{tab:routing_scoring_ablation}
\end{table}

\section{Hyperparameter Optimization}

Hyperparameters are configuration choices that are not learned directly by the
model, but strongly affect training and inference. Examples include the learning
rate, number of epochs, hidden size, class weights, thresholds, split strategy,
and decoder parameters. In this project, hyperparameter optimization was done in
a controlled and experimental way rather than as a large automatic search.

This choice was deliberate. The available learner-error data is limited, so a
large search can easily overfit the validation set. For that reason, the project
used targeted tuning driven by error analysis: first inspect the failure mode,
then tune the parameter that directly affects it, then compare against the
current baseline.

\begin{table}[H]
\centering
\small
\begin{tabularx}{\textwidth}{p{0.24\textwidth} p{0.24\textwidth} p{0.24\textwidth} X}
\toprule
\textbf{Hyperparameter} & \textbf{What was varied} & \textbf{Observed effect} & \textbf{Final interpretation} \\
\midrule
Learning rate &
Especially low learning rate for content auxiliary tuning, for example
$10^{-5}$. &
Aggressive fine-tuning damaged chunk reliability, while low learning rate
preserved the baseline. &
Small learning rates are safer when adapting a content model with limited data. \\
\midrule
Number of epochs &
Short auxiliary training versus longer aggressive training. &
The successful tiny auxiliary content run used only one epoch and still improved
exact match from 0.700 to 0.707. &
More epochs are not automatically better; they can overfit small auxiliary
datasets. \\
\midrule
Decoder blank penalty &
Compared blank penalty 1.0 and 1.6 for chunked CTC decoding. &
Exact match improved from 0.691 to 0.700 and edit distance improved from 0.741
to 0.731. &
The blank penalty is important because CTC models can delete too many
characters. \\
\midrule
Decision thresholds &
Transition, localized support, and learned routing thresholds were tuned. &
Learned router thresholds improved weak-label exact match from 0.895 to 0.921;
multi-label transition used ikhfa 0.45 and idgham 0.40 experimentally. &
Threshold tuning is useful, but a tuned model is promoted only if it improves
strict validation. \\
\midrule
Class weights &
Used to reduce imbalance between frequent and rare labels. &
Transition benefited from class-aware training because idgham had fewer samples
than none and ikhfa. &
Class weights help minority rules, but they must be checked against overall
accuracy. \\
\midrule
Hardcase weights &
Boosted examples that appeared in repeated high-confidence mistakes. &
Transition improved strongly, while duration hardcase variants showed that too
much weighting can hurt balance. &
Hardcase weighting is useful only with a promotion gate; it is not automatically
safe. \\
\midrule
Model size &
Compared smaller and slightly larger content configurations, including the HD96
chunked model. &
The HD96 chunked model became the official content baseline with exact 0.700
and character accuracy 0.892. &
Increasing capacity helps only when paired with strict text-held-out evaluation. \\
\midrule
Data split strategy &
Used text-held-out and verse-held-out splits instead of only random splits. &
Duration fusion was accepted only after verse-held-out accuracy improved from
0.954 to 0.973. &
Strict splits are essential because random splits can hide memorization. \\
\bottomrule
\end{tabularx}
\caption{Hyperparameters tuned during the project and their practical effect.}
\label{tab:hyperparameter_summary}
\end{table}

In the final report, the safest way to describe this work is the following:
the project performed \textit{targeted hyperparameter tuning with controlled
promotion gates}. This means that hyperparameters were not changed randomly.
They were adjusted to solve observed problems, and the new version was accepted
only when it improved the relevant validation metric.

An Optuna-based automatic search is a strong future extension. Optuna could
search over learning rate, thresholds, class weights, blank penalty, hidden
dimension, and loss weights. However, it should be used only after creating
larger and stricter validation sets, otherwise it may optimize the validation
set without improving real generalization.

\section{Comparison with Related Work}

Compared with many existing Quran recitation projects, the proposed system is
more focused on modular diagnosis than on producing a single global correctness
score. Some works focus on automatic speech recognition, some focus on a small
set of Tajweed rules, and others focus on clean recitation. This project
combines several components:

\begin{itemize}[leftmargin=*]
    \item learner-oriented data through RetaSy,
    \item Quran JSON and colored Mushaf references for rule positions,
    \item MFCC-based rule modules,
    \item Wav2Vec-style content recognition,
    \item localized evidence models,
    \item hardcase mining,
    \item severity-aware feedback,
    \item batch evaluation and JSON analysis reports.
\end{itemize}

The main advantage is not only the final metric value, but the fact that each
part can be evaluated separately. This makes the system easier to explain,
debug, extend, and defend in an academic setting.

\section{BERTScore and Confidence Metrics}

BERTScore is useful for semantic similarity in natural language generation.
However, Quran recitation assessment is not mainly a semantic task. A small
letter mistake can be very important even if the semantic embedding similarity
remains high. Therefore, BERTScore is not the main metric for this project.

Instead, the system uses metrics that are closer to the task:

\begin{itemize}[leftmargin=*]
    \item exact match for content,
    \item character accuracy,
    \item edit distance,
    \item class accuracy for rules,
    \item confidence scores from model probabilities,
    \item severity-weighted error penalties.
\end{itemize}

These metrics better reflect whether the learner recited the expected text and
whether the expected Tajweed rule was detected.

\section{Scalability}

The current system is a research prototype. It can scale in three directions.

First, data scalability requires more distinct verses and more learner mistake
examples. The full Quran contains more than six thousand verses, while the
current learner-oriented training data covers a much smaller subset. More ayahs
and more reciters would improve generalization.

Second, model scalability requires moving content recognition from chunked
recognition toward less constrained learner-level recognition. The chunked
baseline is useful now, but the long-term goal is to recognize and score longer
open recitations.

Third, system scalability requires keeping the modular architecture clean. Each
module should have its own training script, evaluation script, manifests,
checkpoints, and analysis reports. This is why the codebase was reorganized into
module-oriented folders.

\section{Final Discussion}

The current system is not a finished commercial product, but it is a complete
development and experimentation prototype. It implements the architecture
described in the conceptual framework: routing, specialist modules, feature
extraction, temporal modeling, diagnosis, feedback generation, and batch
evaluation.

The strongest validated parts are the duration, transition, and burst modules.
The content module improved significantly after moving from full-verse
recognition to chunked recognition, but open learner-level content recognition
remains the main future challenge. The most important scientific lesson of the
project is that larger clean recitation data is useful, but learner mistake data
is essential for true Tajweed assessment.

In summary, the project demonstrates a modular and extensible architecture for
automatic Tajweed assessment. It provides measurable results, controlled
ablations, documented negative experiments, and a clear path for future
improvement.

% -- Bibliography --------------------------------------------
\cleardoublepage
\addcontentsline{toc}{chapter}{Bibliography}
\bibliography{bibliography}

\end{document}
